%!TEX root = ../main/main.tex

%%%%%%%%%%%%%%%%%%%%%%%%%%%%%%%%
%			PACKAGES
%%%%%%%%%%%%%%%%%%%%%%%%%%%%%%%%

\usepackage[utf8]{inputenc}
%\usepackage[english]{babel}
\usepackage{xspace}
\usepackage{algorithm}
\usepackage[noend]{algpseudocode}
\usepackage{amsmath}
\usepackage{amsthm}
\usepackage{amsfonts}
\usepackage{MnSymbol}
\usepackage{mathtools}
\usepackage{textcomp}
\usepackage{stmaryrd}
\usepackage{varwidth}
\usepackage{graphicx}
\usepackage{mathrsfs} 

\usepackage{tikz}
\usetikzlibrary{automata, graphs,positioning,chains,arrows,snakes,decorations.pathmorphing}

%%%% TO CENTER THE EQUATIONS IN LIPICS
\makeatletter
\@fleqnfalse
\@mathmargin\@centering
\makeatother


%%%%%%%%%%%%%%%%%%%%%%%%%%%%%%%%
%			COMMENTS
%%%%%%%%%%%%%%%%%%%%%%%%%%%%%%%%

\usepackage{todonotes}
\newcommand{\Fernando}[1]{\todo[inline, color=violet!20]{{\bf Fernando:} #1}}
\newcommand{\cristian}[1]{\todo[inline, color=red!20]{{\bf Cristian:} #1}}

%IF YOU WANT TO REMOVE COMMENTS, UNCOMMENT THE NEXT LINES
%\renewcommand{\martin}[1]{}
%\renewcommand{\cristian}[1]{}


%%%%%%%%%%%%%%%%%%%%%%%%%%%%%%%%
%			SETS
%%%%%%%%%%%%%%%%%%%%%%%%%%%%%%%%

\newcommand{\cL}{\mathcal{L}}
\newcommand{\cO}{\mathcal{O}}
\newcommand{\cD}{\mathcal{D}}

\newcommand{\bbN}{\mathbb{N}}
\newcommand{\bbR}{\mathbb{R}}
\newcommand{\bbZ}{\mathbb{Z}}



%%%%%%%%%%%%%%%%%%%%%%%%%%%%%%%%
%			ENVIROMENTS
%%%%%%%%%%%%%%%%%%%%%%%%%%%%%%%%

\renewcommand{\paragraph}[1]{\smallskip
	\noindent\textbf{#1}.}
	
%%%%%%%%%%%%%%%%%%%%%%%%%%%%%%%%
%			Streams, CER 
%%%%%%%%%%%%%%%%%%%%%%%%%%%%%%%%

\newcommand{\cA}{\mathcal{A}}

\newcommand{\Stream}{\bar{S}}
\newcommand{\cstart}{\operatorname{start}}
\newcommand{\cend}{\operatorname{end}}
\newcommand{\cinterval}{\operatorname{interval}}
\newcommand{\Pbasic}{\mathbf{P}_{\operatorname{basic}}}
\newcommand{\att}{\operatorname{att}}
\newcommand{\type}{\operatorname{type}}
\newcommand{\ce}{\operatorname{ce}}

\newcommand{\SETT}{\mathbf{T}}
\newcommand{\SETA}{\mathbf{A}}
\newcommand{\SETD}{\mathbf{D}}
\newcommand{\SETX}{\mathbf{X}}
\newcommand{\SETE}{\mathbf{E}}
\newcommand{\SETP}{\mathbf{P}}
\newcommand{\SETL}{\mathbf{L}}

\newcommand{\sem}[1]{\text{\textlbrackdbl} #1 \text{\textrbrackdbl}}

\newcommand{\overlapsimb}{\text{:o}}
\newcommand{\overlap}{~\overlapsimb~}
\newcommand{\cAoverlap}{\cA_{\overlapsimb}}
\newcommand{\Qoverlap}{Q_{\overlapsimb}}
\newcommand{\Poverlap}{\SETP_{\overlapsimb}}
\newcommand{\Xoverlap}{\SETX_{\overlapsimb}}
\newcommand{\Doverlap}{\Delta_{\overlapsimb}}	
\newcommand{\celplus}[1]{\text{CEL} \cup \{#1\}}

\newcommand{\startssimb}{\text{:s}}
\newcommand{\starts}{~\startssimb~}
\newcommand{\cAstarts}{\cA_{\startssimb}}
\newcommand{\Qstarts}{Q_{\startssimb}}
\newcommand{\Pstarts}{\SETP_{\startssimb}}
\newcommand{\Xstarts}{\SETX_{\startssimb}}
\newcommand{\Dstarts}{\Delta_{\startssimb}}	

%%%%%%%%%%%%%%%%%%%%%%%%%%%%%%%%
%			ALGORITHMS
%%%%%%%%%%%%%%%%%%%%%%%%%%%%%%%%

\renewcommand{\operatorname}[1]{\mathsf{#1}}

\algrenewcommand\algorithmicindent{1.3em}%
\algnewcommand\algorithmicforeach{\textbf{for each}}
\algdef{S}[FOR]{ForEach}[1]{\algorithmicforeach\ #1\ \algorithmicdo}



